%=========== ΛΥΣΗ - 93ο ΘΕΜΑ Β ===========

\begin{thema}{Β}
\begin{erwthma}
\item Από το σχήμα της $f''$ βλέπουμε ότι η γραφική παράσταση τέμνει τον άξονα $x'x$ στο $x=1$. Επίσης βρίσκεται κάτω από τον άξονα για $x<1$ και πάνω από τον άξονα για $x>1$. Άρα
\[
f''(x)<0 \text{ για } x<1,\quad f''(1)=0,\quad f''(x)>0 \text{ για } x>1.
\]
Επομένως η $f$ είναι κοίλη στο $(-\infty,1]$ και κυρτή στο $[1,+\infty)$. Επειδή στο $x=1$ αλλάζει η κυρτότητα, η $C_f$ έχει σημείο καμπής το $K(1,f(1))$.

\item Η $f$ παρουσιάζει τοπικό ακρότατο στο $x_1=0$ και είναι παραγωγίσιμη, άρα από το θεώρημα \eng{Fermat} ισχύει $f'(0)=0$. Από το σχήμα διαβάζουμε ότι στο $x=0$ η τιμή της δεύτερης παραγώγου είναι αρνητική, δηλαδή $f''(0)<0$. Επομένως, με το κριτήριο της δεύτερης παραγώγου, η $f$ παρουσιάζει τοπικό μέγιστο στο $x=0$.

\item Από το προηγούμενο ερώτημα $f'(0)=0$ και δίνεται $f'(2)=0$. Άρα, από το Θεμελιώδες Θεώρημα,
\[
\int_0^2 f''(x)\,\mathrm{d}x=\left[f'(x)\right]_0^2=f'(2)-f'(0)=0.
\]
\end{erwthma}
\end{thema}

%=========== ΤΕΛΟΣ ΛΥΣΗΣ - 93ο ΘΕΜΑ Β ===========
